\section{RREF}

The matrix on the handout was printed as a stacked list of nine entries. We
interpret it as the $3\times 3$ matrix obtained by reading those entries in
row-major order
\[
A=\begin{pmatrix}0&1&0\\3&0&1\\6&5&2\end{pmatrix}.
\]
(The algorithm and structural theorems below do not depend on this choice.)

\subsection{Q1. The RREF algorithm, and $\RREF(A)$}

\begin{promptbox}
State the RREF algorithm and compute the RREF of
$A=\begin{pmatrix}0&1&0\\3&0&1\\6&5&2\end{pmatrix}$.
\end{promptbox}

\begin{coursebox}
\textbf{Algorithm} (Gaussian elimination to reduced row echelon form).
Start with $A$ and a ``current'' top-left position $(r,c)=(1,1)$. Repeat until
you run out of rows or columns:
\begin{enumerate}[leftmargin=1.8em]
\item Among rows $r,\ldots,m$, find a row whose entry in column $c$ is nonzero.
If none exists, increment $c$ and repeat.
\item Swap that row with row $r$ (a pivot now sits at $(r,c)$).
\item Scale row $r$ so the pivot equals $1$.
\item Add multiples of row $r$ to every other row so that all other entries in
column $c$ become $0$.
\item Increment $r$ and $c$ and continue.
\end{enumerate}
The output is $\RREF(A)$: each pivot is $1$, each pivot is the first nonzero
entry of its row, and each pivot column is a standard basis vector.

\textbf{Computation.}
Column $1$ is nonzero. Swap row $1$ with row $2$, then scale the new row $1$ by
$1/3$:
\[
\begin{pmatrix}3&0&1\\0&1&0\\6&5&2\end{pmatrix}
\to
\begin{pmatrix}1&0&1/3\\0&1&0\\6&5&2\end{pmatrix}.
\]
Subtract $6$ times row $1$ from row $3$:
\[
\begin{pmatrix}1&0&1/3\\0&1&0\\0&5&0\end{pmatrix}.
\]
Subtract $5$ times row $2$ from row $3$:
\[
\begin{pmatrix}1&0&1/3\\0&1&0\\0&0&0\end{pmatrix}.
\]
Pivots are already $1$ and already clear above. So
\[
\RREF(A)=\begin{pmatrix}1&0&1/3\\0&1&0\\0&0&0\end{pmatrix}.
\]
\end{coursebox}

\subsection{Q2. $\RREF(A)=GA$ for some invertible $G$}

\begin{promptbox}
Prove that $\RREF(A)=GA$ for some invertible $G$.
\end{promptbox}

\begin{coursebox}
Each elementary row operation is left-multiplication by an elementary matrix
$E$ (swap: a transposition matrix; scale row $i$ by $\lambda\neq 0$: the
diagonal matrix with $\lambda$ in position $i$; add $\lambda$ times row $j$ to
row $i$: $I+\lambda E_{ij}$). Every elementary matrix is invertible (the inverse
operation is elementary). The algorithm applies finitely many such operations,
so $\RREF(A)=E_k\cdots E_1 A=GA$ with $G=E_k\cdots E_1$ invertible.
\end{coursebox}

\subsection{Q3. $\ker(T_A)=\ker(T_{\RREF(A)})$}

\begin{promptbox}
Prove that $\ker(T_A)=\ker(T_{\RREF(A)})$.
\end{promptbox}

\begin{coursebox}
By Q2, $\RREF(A)=GA$ with $G$ invertible. Then $A\vec{x}=\zero$ iff
$GA\vec{x}=\zero$ (multiply by $G$ or by $G^{-1}$). Thus
$\ker T_A=\ker T_{GA}=\ker T_{\RREF(A)}$.
\end{coursebox}

\subsection{Q4. Example with $\im(T_A)\neq\im(T_{\RREF(A)})$}

\begin{promptbox}
Find an example of $A$ with $\im(T_A)\neq\im(T_{\RREF(A)})$.
\end{promptbox}

\begin{coursebox}
Row operations change column space in general. Take
\[
A=\begin{pmatrix}1&0\\1&0\end{pmatrix},\qquad
\RREF(A)=\begin{pmatrix}1&0\\0&0\end{pmatrix}.
\]
Then $\im(T_A)=\operatorname{Span}\{(1,1)^T\}$ while
$\im(T_{\RREF(A)})=\operatorname{Span}\{(1,0)^T\}$, which are unequal.
\end{coursebox}

\subsection{Q5. Same image without $A$ already in RREF}

\begin{promptbox}
Is it possible that $\im(T_A)=\im(T_{\RREF(A)})$ but $A\neq\RREF(A)$?
\end{promptbox}

\begin{coursebox}
Yes. Take
\[
A=\begin{pmatrix}2&0\\0&0\end{pmatrix},\qquad
\RREF(A)=\begin{pmatrix}1&0\\0&0\end{pmatrix}.
\]
The column spaces are both $\operatorname{Span}\{(1,0)^T\}$, but $A\neq\RREF(A)$.

(The phenomenon is: scaling a pivot column of $A$ does not change its span, but
does change the matrix.)
\end{coursebox}

\subsection{Q6. Row-equivalent matrices have the same RREF}

\begin{promptbox}
True/false: If $A=GB$ for some invertible $G$, then $\RREF(A)=\RREF(B)$.
\end{promptbox}

\begin{coursebox}
\textbf{True.}
$A=GB$ with $G$ invertible means $A$ and $B$ are row-equivalent: they have the
same row space and, more to the point, each is obtained from the other by
elementary row operations. The reduced row echelon form of a matrix is
\emph{unique}, so there is only one RREF in a row-equivalence class. Hence
$\RREF(A)=\RREF(B)$.

(If uniqueness of RREF has not been proved yet: $\RREF(A)=G_A A=G_A G B$ and
$\RREF(B)=G_B B$ are both RREF matrices that are row-equivalent, and one checks
that a matrix in RREF is determined by its pivot pattern and row space, or
invokes uniqueness.)
\end{coursebox}

\begin{laterbox}
Uniqueness of RREF is the key lemma. A quick uniqueness argument: if $R$ and $S$
are both RREF and $R=HS$ for invertible $H$, then $R$ and $S$ have the same
nullspace, hence the same pivot columns (the non-free coordinates), and comparing
pivot columns --- which are standard basis vectors --- forces $H=I$ on the
pivot rows, hence $R=S$.
\end{laterbox}

\subsection{Q7. RREF when $T_A$ is injective}

\begin{promptbox}
If $T_A$ is injective, what does $\RREF(A)$ look like?
\end{promptbox}

\begin{coursebox}
$T_A:\F^n\to\F^m$ injective means $\ker T_A=\{\zero\}$, so there are no free
variables: every column is a pivot column. That requires $n\le m$, and
\[
\RREF(A)=\begin{pmatrix}I_n\\0_{(m-n)\times n}\end{pmatrix}
\]
(the block $0$ is absent if $m=n$, in which case $\RREF(A)=I_n$).
\end{coursebox}

\subsection{Q8. Bases of $\im(T_A)$ and $\ker(T_A)$ from $\RREF(A)$}

\begin{promptbox}
Given $\RREF(A)$, describe how to find a basis for $\im(T_A)$ and $\ker(T_A)$.
\end{promptbox}

\begin{coursebox}
Let $R=\RREF(A)$.
\begin{itemize}
\item \textbf{Image / column space of $A$.}
The pivot columns of $R$ mark the pivot columns of $A$. Those columns of the
\emph{original} matrix $A$ form a basis of $\im(T_A)$. (Do not take columns of
$R$: Q4 shows they may span the wrong space.)
\item \textbf{Kernel / nullspace of $A$.}
Free variables correspond to non-pivot columns. For each free variable $x_j$,
set $x_j=1$ and set the other free variables to $0$; solve the pivot equations
(which are immediate from $R$) for the pivot variables. The resulting $n$
special solutions form a basis of $\ker(T_A)$. Equivalently, if the non-pivot
columns of $R$ are $R\vec{e}_{j_k}$, each special solution is
$\vec{e}_{j_k}$ minus the linear combination of pivot columns that $R$ uses to
express that column.
\end{itemize}

\textbf{Illustration with the matrix of Q1.}
$R=\begin{pmatrix}1&0&1/3\\0&1&0\\0&0&0\end{pmatrix}$ has pivots in columns
$1,2$. A basis of $\im(T_A)$ is the first two columns of $A$:
$(0,3,6)^T$ and $(1,0,5)^T$. One free variable $x_3$, and
$x_1+(1/3)x_3=0$, $x_2=0$, so a basis of the kernel is
$\bigl((-1/3),\,0,\,1\bigr)^T$.
\end{coursebox}
